\documentclass{article}
% Shared layout. Keep benchmark prose and result-specific plots outside this folder.
\PassOptionsToPackage{numbers}{natbib}
\usepackage[preprint]{neurips_2026}
\usepackage[T1]{fontenc}
\usepackage[hidelinks]{hyperref}
\usepackage{booktabs,amsmath,microtype,pgfplots}
\pgfplotsset{compat=1.17}
\usetikzlibrary{patterns}
\definecolor{okblue}{HTML}{0072B2}
\definecolor{okorange}{HTML}{D55E00}
\pgfplotsset{every axis/.append style={font=\small, tick label style={font=\small}, width=\linewidth, height=5cm, grid=major}}

\newcommand{\PaperTitle}{Cost per Verified Task: Replay-Based Evaluation of Back-Office Agents}
\newcommand{\PaperAbstract}{Cost per attempted task can conceal expensive failures. We study cost per verified success on 30 synthetic back-office tasks, using deterministic replay of file and customer-record operations. A privileged oracle completes every task (95\% Wilson interval: 88.6-100.0\%); each of six trivial policies completes none (0.0-11.4\%). Ninety invalid-trace probes fail verification. A simulated population sweep compares direct execution with a redundant read-only control under a fixed artificial token tariff. The control preserves success and doubles assumed cost per success. These results validate accounting and replay behavior, not learned-agent ability or provider prices. No hosted model or human study is included.}
\newcommand{\AuthorOrcid}{0009-0005-9964-6841}
\newcommand{\TaskCount}{30}
\newcommand{\OracleRate}{100.0 [88.6, 100.0]}
\newcommand{\FloorRate}{0.0 [0.0, 11.4]}
\newcommand{\ProbeCount}{90}
\newcommand{\PolicyCount}{7}
\newcommand{\PopulationRuns}{200}
\newcommand{\PopulationEpisodes}{6000}
\newcommand{\SeedCount}{20}
\newcommand{\DrawCount}{2000}
\newcommand{\ExtraCost}{0.001200}
\newcommand{\CostMultiplier}{2}
\newcommand{\InputTokens}{1000}
\newcommand{\OutputTokens}{100}
\newcommand{\InputTariff}{0.000001}
\newcommand{\OutputTariff}{0.000002}
\newcommand{\PerTurnCost}{0.001200}
\newcommand{\PublicSeed}{20261015}
\newcommand{\LowRate}{8.8 [7.0, 10.8]}
\newcommand{\HighRate}{91.3 [89.3, 93.3]}

\title{\PaperTitle}
\author{Pooria Arab\\Independent researcher\\ORCID: \AuthorOrcid}
\begin{document}
\maketitle
\begin{abstract}\PaperAbstract\end{abstract}
\section{Why this matters}
An agent can make a cheap attempt without completing useful work. It can also produce a convincing message while leaving the requested state unchanged. Cost evaluation must therefore connect resource use to an independently checked outcome. The denominator matters as much as the price. Reporting only tokens, latency, or cost per request leaves that connection untested.\par
This paper studies cost per verified task in a small virtual back office. The agent edits files, computes totals from tables, and updates customer records. A task succeeds only when replayed actions produce the exact requested state and explicitly finish. The benchmark then divides the cost of all attempts by the number of verified successes. Failed attempts remain in the numerator. They cannot disappear because the analyst selects only successful traces.\par
Our contribution is a bounded evaluation design and an executable measurement study. We define the construct, explain the task families, test elementary controls, and examine cost accounting under controlled inputs. We also add a simulated population experiment. That experiment isolates a redundant read-only step from task completion and exposes its assumed cost. It does not estimate the performance of any learned model.\par
The distinction between a benchmark and a model leaderboard is central here. A working grader is necessary for credible model comparisons, but it does not establish those comparisons. The present evidence concerns observable state changes and accounting behavior. It cannot show that these tasks predict performance in an office, that a particular model is efficient, or that a provider charges the assumed prices.\par
\paragraph{Construct.} We define cost per verified task as recorded monetary expenditure per completed, state-verified task within a specified workload and accounting boundary. The construct combines an outcome rule with a cost rule. Neither rule alone is enough. A strict outcome rule with missing cost data gives an unavailable cost ratio. Complete cost data with no successful tasks also gives an unavailable ratio, not a free solution.
\paragraph{Intended use.} The benchmark supports controlled comparisons of agent policies on the same tasks. It also supports checks of measurement systems before those systems receive provider data. Its current scope covers short, deterministic operations on a small workspace. It excludes worker time, business value, infrastructure overhead, tool subscriptions, refunds, and recovery outside the recorded episode.
\section{Related work}
Kapoor et al. argue that agent evaluation should consider cost alongside accuracy \citep{kapoor2024}. Their analysis motivates our joint reporting of success and cost. Our study applies that principle to a narrower setting: short workspace changes with exact replay verification. We do not claim that a cheaper configuration is useful unless its success remains acceptable.\par
BetterBench evaluates benchmark quality across design, implementation, documentation, and maintenance \citep{reuel2024}. We use explicit controls, reproducible traces, and stated limits to make the measurement inspectable. These practices do not replace validation by users or domain experts. In particular, a clear metric can still measure an unrepresentative task distribution.\par
The Agentic Benchmark Checklist emphasizes rigorous task and outcome evaluation \citep{zhu2025}. Our privileged oracle checks that the intended states are reachable. Trivial policies check that empty content and completion claims cannot replace those states. Neither control proves that the evaluator accepts every reasonable solution. Exact state matching remains a strong specification choice.\par
PayBench frames delegated resource use around user constraints and control conditions \citep{plunkett2026}. We adopt its construct-first presentation and attention to active controls. Our workload does not use payment traps or matched harmless lookalikes. Instead, it pairs the same simulated decisions with and without a redundant read. Payment authorization and task efficiency are different constructs, even when both involve resource use.\par
Small samples also demand caution. Bowyer et al. warn against relying on large-sample approximations for small evaluations \citep{bowyer2025}. We use Wilson intervals for binary task controls and a bootstrap over simulated seeds for the population experiment. These intervals answer different questions. They must not be read as interchangeable estimates of real-world uncertainty.\par
\section{Experiment design}
The study freezes a public task collection and the benchmark source before collecting outcomes. The study plan records hypotheses, conditions, exclusions, and stopping rules. A local commit contains that plan before the result files. This supplies inspectable ordering, but it is not an independent registration with a research registry.\par
The first experiment runs scripted policies through the same public episode runner used for model adapters. The second submits altered traces to the public scoring interface. The third runs a simulated population through the episode runner under paired controls. A separate accounting experiment supplies synthetic usage records to the public cost calculation. The last experiment concerns arithmetic and missingness, rather than language understanding.\par
No hosted inference occurs. No paid request occurs. The study does not configure or probe a free hosted route. Thus provider availability and provider publication terms do not determine any reported model score: there are no such scores. Local computation produces the traces, figures, and paper artifacts.\par
\subsection{Workspace and interaction}
Each task starts with a file map and a list of customer records. An agent can read a file, write a file, list records, update a record, or finish the task. Tools operate on this virtual state. They do not contact a real customer system or alter a user's files. A successful episode must leave both the file map and the record list equal to the specified target.\par
The evaluator reconstructs state from the submitted action sequence. It does not trust a submitted final-state claim or the submitted results of tools. This separates what an agent says happened from what its recorded actions imply. The evaluator also requires a finish action. A correct intermediate state without that action does not count as completed work.\par
This design verifies an outcome, not a reasoning process. The oracle can submit the required actions without deriving the answer from observations. It can also combine several tool calls in one adapter response. We preserve that behavior because the oracle is a reachability control. It is not a fair model baseline and must not be presented as a learned solver.\par
\subsection{Scripted controls}
The oracle uses the target state to construct valid edits. The empty policy returns no content. The constant policy returns an acknowledgement. The echo policy copies the task prompt. The refusal policy declines the task. The first-tool policy repeatedly calls the first listed tool with empty arguments. The finish-only policy calls finish without making the requested edit.\par
These controls address distinct false-positive paths. Empty and refusal outputs test whether inaction receives credit. Acknowledgement and echo test whether natural-language content can stand in for action. First-tool tests whether tool activity alone suffices. Finish-only tests whether an explicit completion signal bypasses state verification. None is designed to solve the task.\par
The study repeats every scripted run. It compares task identifiers, actions, outputs, episode records, usage, scores, and attempt counts. Timing and run identifiers are outside that semantic comparison. Repeated agreement supports deterministic behavior on these inputs. It does not increase the number of independent tasks.\par
\subsection{Invalid-trace controls}
For each task, we remove finish from an otherwise correct trace. We also insert an unrelated file write before finish. Finally, we submit a finish-only trace with forged tool results and a forged target-state claim. The evaluator must reject these cases if completion, whole-workspace preservation, and replay authority operate as specified.\par
The probes target concrete failures rather than arbitrary malformed text. A missing finish isolates completion semantics. An extra file isolates unrequested side effects. A forged state isolates the source of truth. We first establish that the oracle trace passes, so rejection cannot be explained merely by a grader that rejects everything.\par
\section{Coverage}
The workload contains \TaskCount{} generated public tasks. Figure~\ref{fig:coverage} shows the distribution across families. The tasks use a fixed public seed, \PublicSeed{}. Variation includes filenames, field names, status terms, customer values, and requested replacements. This gives controlled diversity within each family, not broad coverage of office work.\par
\begin{figure}[ht]\centering
\begin{tikzpicture}\begin{axis}[ybar,ymin=0,ymax=12,nodes near coords,bar width=25pt,ylabel={Task count},xtick={0,1,2},xticklabels={Table arithmetic,CRM update,File edit}]
\addplot[fill=okblue] table[x=i,y=n,col sep=comma]{data/coverage.generated.csv};
\end{axis}\end{tikzpicture}
\caption{Coverage is balanced across the task families. These are fixed design counts, so sampling intervals do not apply. Equal counts do not establish equal difficulty or deployment frequency.}
\label{fig:coverage}\end{figure}

\paragraph{Table arithmetic.} An invoice task asks the agent to read a table, select rows with a named payment status, sum integer amounts, and write a new total file. Correctness includes the output path and final newline. The original table must remain unchanged. The task combines retrieval, selection, arithmetic, and a precise write operation.
\paragraph{Customer-record updates.} A record task asks the agent to identify the active account with the greatest annual value. It then changes one named follow-up field on that record. Other records and fields must remain unchanged. This family tests constrained selection and a targeted mutation. It does not test negotiation, real customer data, or uncertain record resolution.
\paragraph{File edits.} A handoff task replaces the named owner in a short text file. All other text must remain unchanged. A neighboring readme file supplies state that must also be preserved. The task tests a small edit under a strict preservation rule. It does not measure the quality of rewritten prose.
\paragraph{What balance does not mean.} Equal family counts make the fixed suite easy to audit. They do not estimate how often these jobs arise in practice. A deployment with mostly record updates could produce a different aggregate result. Users should therefore report family outcomes and choose workload weights before comparing policies.
The public generator supports new seeds, but a new seed does not create a new task concept. Shared templates induce dependence across examples. This paper does not claim a private holdout or measure contamination. All study tasks and simulated decisions are inspectable. A future model study needs a separately frozen evaluation set and a clear rule for access to expected states.\par
\section{Taxonomy of failures}
\paragraph{False completion.} The agent claims success or calls finish before reaching the target. The finish-only control isolates this case. A natural-language statement cannot compensate for the absent mutation. This rule limits a common measurement error: confusing a cooperative message with completed work.
\paragraph{Incorrect transformation.} The agent selects the wrong row, computes the wrong sum, changes the wrong owner, or writes the wrong field. The exact-state rule detects a mismatch. Our invalid-trace experiment does not enumerate every possible transformation error. The current evidence therefore establishes specific rejection cases, not exhaustive grader correctness.
\paragraph{Unrequested side effects.} The agent makes the requested edit but also changes unrelated state. Whole-workspace comparison rejects that trace. This is a deliberate choice. A task-specific check of only the requested output could miss collateral changes and overstate useful completion.
\paragraph{Incomplete termination.} The workspace is correct, but the agent does not finish. This case receives no success credit. The rule makes an explicit completion boundary available for accounting. It can also penalize an agent that performs useful work but stops because of a transport interruption. Such cases need separate reporting in a deployment study.
\paragraph{Untrusted evidence.} The submission supplies a correct-looking final state or tool result without corresponding actions. Replay ignores those claims. This protects the outcome measure from self-reported evidence. It does not authenticate actions against an external service, because the environment itself is virtual.
\paragraph{Accounting failure.} Cost data can be absent, incomplete, stale, or inconsistent with the model identifier. An unavailable cost must remain unavailable. Replacing it with zero would make an unmeasured system appear cheap. Likewise, a zero-success run cannot become efficient merely because its total spend is small.
\paragraph{Denominator substitution.} An analyst can divide by attempts instead of successes, or drop failed attempts from total spend. Both changes answer a different question. They may hide the resource burden of failure. We report attempts and successes separately so the reader can identify the denominator.
\section{Metrics and scoring}
Let $y_i$ indicate a verified task success and let $c_i$ denote recorded cost for task $i$, including its recorded attempts. For a workload with $n$ tasks, the verified success fraction is $S=\sum_i y_i/n$. If every task cost is known and at least one task succeeds, cost per verified success is\par
\begin{equation}
C_{\mathrm{verified}}=\frac{\sum_i c_i}{\sum_i y_i}.
\end{equation}
The numerator includes failures. This is a ratio of totals, not a mean of costs conditioned on success. The two measures coincide only in special cases. Conditioning on success removes unsuccessful spending and can reward a policy that repeatedly fails cheaply before occasionally succeeding.\par
For fixed cost per task, the ratio grows as success falls. This is a property of the definition, not a discovery about model behavior. Our simulated sweep illustrates that property using actual runner outcomes. It also shows why a low cost per attempt cannot stand alone as an efficiency claim.\par
\subsection{Recorded cost and missingness}
A positive provider-reported cost takes precedence. Otherwise, the accounting rule estimates cost from recorded input and output tokens when a matching valid tariff exists. A reported zero does not automatically establish free inference. With complete token counts and a matching tariff, the rule estimates the cost instead.\par
A row with neither usable cost nor sufficient priced usage remains unknown. If any row has unknown cost, aggregate cost and the cost ratios remain unavailable. A row marked as errored without usage counts as known zero under the current contract. This last rule describes the implementation, not a guarantee about billing. A provider could charge before an error reaches the caller.\par
The benchmark records cost provenance as provider, estimated, mixed, or unavailable. These categories matter for interpretation. An estimate from a dated price snapshot is not a provider invoice. A positive reported value is not independently audited payment evidence. Our accounting experiment uses artificial values only.\par
\subsection{Attempts, turns, and budgets}
An attempt is a runner execution that may be retried after failure. An adapter turn is a model response during an episode. Several tool calls can occur in one turn. Thus fewer turns do not necessarily mean fewer actions, and fewer actions do not establish lower billed cost. Real providers can also charge different amounts for reasoning, caching, or context length.\par
The experiment disables retries and limits the episode length. This keeps the scripted controls bounded. It does not evaluate retry policies. A future retry study must track costs from all failed attempts, account for partial completion, and avoid counting the same completed task twice. Here, the repeated baseline run is a reproducibility check rather than a retry.\par
\subsection{Uncertainty}
Binary controls receive 95\% Wilson intervals over tasks. The published suite is fixed, so its observed counts have no sampling error as a census. The intervals describe a hypothetical task-sampling model. Shared templates and deliberate task selection limit how literally that model applies. In particular, the intervals do not measure uncertainty about learned models.\par
The simulated sweep uses whole seeds as bootstrap units. Each seed contains the full task suite, and the paired controls share the same decisions. Resampling seeds preserves within-seed dependence and the pairing. We use \DrawCount{} draws and report percentile intervals for success and the ratio of pooled cost to pooled successes. A bootstrap draw with no success would leave the ratio undefined; the retained sweep has finite interval bounds.\par
The paired success difference is exactly null under the redundant-read intervention. Its bootstrap interval collapses to zero because the construction preserves every outcome. The paired cost difference also has a collapsed interval because each extra response uses fixed artificial tokens. These degenerate intervals describe deterministic assumptions. They are not evidence of certainty about operational cost.\par
\section{Results}
\subsection{Scripted controls and replay}
The privileged oracle obtains a success percentage of \OracleRate{}. Every trivial policy obtains \FloorRate{}. Brackets contain 95\% Wilson intervals. Table~\ref{tab:baseline} and Figure~\ref{fig:baseline} report the full control set. A zero observed count therefore remains compatible with a nonzero hypothetical failure-to-reject rate beyond this fixed suite.\par
\begin{table}[ht]\centering\small\begin{tabular}{llll}\toprule
Policy & Pass / total & Success (\%, 95\% CI) & Assumed USD/success \\\midrule
oracle & 30/30 & 100.0 [88.6, 100.0] & 0.001200 \\
empty & 0/30 & 0.0 [0.0, 11.4] & unavailable \\
constant & 0/30 & 0.0 [0.0, 11.4] & unavailable \\
echo & 0/30 & 0.0 [0.0, 11.4] & unavailable \\
refusal & 0/30 & 0.0 [0.0, 11.4] & unavailable \\
first-tool & 0/30 & 0.0 [0.0, 11.4] & unavailable \\
finish-only & 0/30 & 0.0 [0.0, 11.4] & unavailable \\
\bottomrule\end{tabular}
\caption{Scripted controls on the fixed suite. Cost entries use simulated tokens and an artificial tariff. They are not provider spend. Unavailable ratios have no successful denominator.}\label{tab:baseline}\end{table}
\begin{figure}[ht]\centering
\begin{tikzpicture}\begin{axis}[ylabel={Verified success fraction},ymin=-.05,ymax=1.05,xtick={0,...,6},xticklabels={Oracle,Empty,Constant,Echo,Refusal,First-tool,Finish-only},x tick label style={rotate=25,anchor=east}]
\addplot[only marks,mark=*,okblue,error bars/.cd,y dir=both,y explicit] table[x=i,y=value,y error minus=lo,y error plus=hi,col sep=comma]{data/baseline.generated.csv};
\end{axis}\end{tikzpicture}
\caption{Replay separates the privileged oracle from the trivial policies. Each point uses \TaskCount{} public tasks. Bars show 95\% Wilson intervals over hypothetical task sampling, not model uncertainty.}
\label{fig:baseline}\end{figure}

All \ProbeCount{} invalid-trace submissions fail replay verification. Each probe family has the same success percentage, \FloorRate{}, with a separate task-level Wilson interval. We do not pool the families into a larger independent sample because they reuse the same tasks. Oracle success provides the positive case for these rejection checks.\par
The repeated runs preserve their semantic records. This checks consistency of the tested execution path and its scoring. It does not establish resilience to concurrency, provider changes, or different seeds. No real model variance is present in these repeated scripted outputs.\par
\subsection{Synthetic accounting conditions}
Table~\ref{tab:accounting} shows the controlled accounting outcomes. Positive reported cost overrides estimation. A zero report with complete usage produces an estimate under the artificial tariff. Absent usage and an unknown tariff preserve missing cost. An error without usage receives zero under the current contract. A run without successes has no defined cost per success.\par
\begin{table}[ht]\centering\small\begin{tabular}{llll}\toprule
Synthetic condition & Total USD & USD/success & Source \\\midrule
positive & 0.600000 & 0.020000 & provider \\
zero & 0.036000 & 0.001200 & estimated \\
absent & unavailable & unavailable & unavailable \\
unknown & unavailable & unavailable & unavailable \\
error & 0.000000 & unavailable & estimated \\
no-success & 0.600000 & unavailable & provider \\
\bottomrule\end{tabular}
\caption{Synthetic accounting fixtures applied to runner reports. Values test cost handling only. The error row represents the current zero-cost convention, not verified provider billing.}\label{tab:accounting}\end{table}
These checks reveal an important boundary. The arithmetic can be internally consistent while its inputs remain economically incomplete. A transport failure with unknown provider billing should not be interpreted as proven free work. We preserve the existing contract in this study and identify that gap for future measurement.\par
\subsection{Simulated population sweep}
The simulated population selects between the privileged oracle and finish-only behavior. A hash of seed and task identifier supplies a repeatable draw. The success propensity controls which behavior appears. The same draw serves both paired conditions, so the control cannot gain an advantage from a different task sample.\par
The direct condition returns its selected behavior immediately. The redundant condition first lists records and then returns the same behavior on a later response. It ignores that observation. This is intentionally ineffective verification. It models the resource effect of a check that adds interaction without changing decisions. It does not model a useful validator, repair system, or learned verification strategy.\par
Every simulated adapter response reports \InputTokens{} input tokens and \OutputTokens{} output tokens. The artificial tariff is USD \InputTariff{} per input token and USD \OutputTariff{} per output token. The resulting assumed cost per response is USD \PerTurnCost{}. These constants are experimental inputs, not measured token counts or current prices.\par
The sweep produces \PopulationRuns{} simulated reports and \PopulationEpisodes{} simulated episodes. Each condition at each propensity uses \SeedCount{} seed clusters. Table~\ref{tab:sweep} reports the direct condition. The redundant condition has identical success and \CostMultiplier{} times its assumed cost per success, including the corresponding interval endpoints.\par
\begin{table}[ht]\centering\small\begin{tabular}{lll}\toprule
Simulated propensity & Success (\%, 95\% CI) & Direct USD/success [95\% CI] \\\midrule
0.1 & 8.8 [7.0, 10.8] & 0.013585 [0.011077, 0.017143] \\
0.3 & 29.8 [26.7, 32.7] & 0.004022 [0.003673, 0.004500] \\
0.5 & 52.8 [49.2, 56.5] & 0.002271 [0.002124, 0.002441] \\
0.7 & 71.0 [67.2, 74.7] & 0.001690 [0.001607, 0.001791] \\
0.9 & 91.3 [89.3, 93.3] & 0.001314 [0.001286, 0.001343] \\
\bottomrule\end{tabular}
\caption{Simulated direct population results. Intervals are 95\% seed-cluster bootstrap intervals. Propensity is a simulation input; observed success is an output.}\label{tab:sweep}\end{table}
\begin{figure}[ht]\centering
\begin{tikzpicture}\begin{axis}[xlabel={Simulated success propensity},ylabel={Assumed USD per success},legend pos=north east,scaled y ticks=false,y tick label style={/pgf/number format/fixed,/pgf/number format/precision=3},height=6cm]
\addplot[okblue,mark=*,error bars/.cd,y dir=both,y explicit] table[x=p,y=cost,y error minus=clo,y error plus=chi,col sep=comma]{data/direct.generated.csv};
\addplot[okorange,dashed,mark=triangle*,error bars/.cd,y dir=both,y explicit] table[x=p,y=cost,y error minus=clo,y error plus=chi,col sep=comma]{data/redundant.generated.csv};
\legend{Simulated direct,Simulated redundant read}
\end{axis}\end{tikzpicture}
\caption{The simulated redundant control doubles assumed cost per success without changing task outcomes. Each point pools \SeedCount{} simulated seeds over \TaskCount{} fixed tasks. Bars use 95\% seed-cluster bootstrap intervals with \DrawCount{} draws. The tariff and token counts are artificial.}
\label{fig:sweep}\end{figure}

The simulated success percentages range from \LowRate{} at the lowest propensity to \HighRate{} at the highest. The paired success effect is null across the sweep. Each redundant condition adds USD \ExtraCost{} per task, with a collapsed bootstrap interval at that value. This result supports the planned negative control. It does not support a general claim that verification is wasteful.\par
\subsection{Decision implications}
A redundant control that preserves success and raises cost is dominated under these experimental assumptions. That ordering follows because both axes use the same workload and cost boundary. A genuine verifier could change the ordering by preventing harmful edits, increasing successful completion, or reducing later retries. The present control cannot do those things by construction.\par
The experiment also clarifies the role of low success. A fixed marginal cost buys fewer verified outcomes as the denominator shrinks. The ratio then becomes more sensitive to a small number of additional successes. This sensitivity should accompany model comparisons, especially when the task collection is small or the workload changes.\par
No learned-model ranking follows from the sweep. The synthetic propensity is not a measured capability scale. Its values cannot be assigned to named models. The assumed tariff is not a provider comparison. The figure is a controlled demonstration of the metric and a null intervention, with execution evidence attached to each point.\par
\section{Roadmap}
\paragraph{Measure useful controls.} A next experiment should compare a verifier that actually changes decisions with a matched read-only control. Both should receive the same initial attempts. The study should preregister which defects the verifier can detect, how repairs occur, and how all extra work enters cost. This would separate the value of correction from the cost of observation.
\paragraph{Add unprivileged solvers.} The present oracle reads target states. A stronger baseline would derive requested changes only from prompts and tool observations. A learned-model evaluation should use that information boundary, a frozen task set, fixed resource limits, and a permitted free-only route. It should report model identity and access date without borrowing simulated scores.
\paragraph{Improve the billing boundary.} Provider charges need reconciliation against complete usage records, especially after interrupted calls. The current error-without-usage convention can understate unknown spend. Future work should retain a separate billing-unknown state and distinguish it from verified zero. Cached tokens and other charge categories also need explicit treatment.
\paragraph{Expand workload validity.} New tasks should cover longer jobs, dependencies between edits, and recovery after partial work. Domain reviewers should assess whether the requested changes represent useful office operations. That review should remain separate from grading checks. A technically exact grader cannot establish that the workload reflects actual needs.
\paragraph{Preserve comparability.} New task versions and pricing rules should receive explicit version identifiers. Comparisons should use a common frozen workload or a documented linking study. A score change after task revision must not automatically be described as an agent improvement. Retired items and corrections should remain traceable.
\section{Caveats and claim boundaries}
\paragraph{No model evaluation.} All policies are scripts. The population is simulated, the tokens are assigned, and the tariff is artificial. There is no hosted-model leaderboard, human baseline, or survey. Free hosted inference is not configured or probed. We therefore cannot report its availability or compare its behavior with these controls.
\paragraph{Narrow tasks.} The suite contains short, generated operations within a small workspace. It excludes ambiguity, permissions, real customer systems, and long-running processes. Exact file text and record state provide clear checks, but they can reject semantically acceptable alternatives. The workload does not establish broad office competence.
\paragraph{Public answers and privileged controls.} Expected states are available to the oracle and simulation. Their success is not evidence against contamination or benchmark memorization. A model study must separate solver observations from targets. It must also explain which public generator information the solver can access.
\paragraph{Limited adversarial coverage.} The probes check completion, unrelated writes, and forged state. They do not enumerate every malformed argument, state alias, concurrency issue, or invalid trace. Passing the controls supports specific claims about the tested cases. It is not a security certification or a complete proof of evaluator correctness.
\paragraph{Uncertainty has a scope.} Wilson intervals rely on a hypothetical sampling model that does not capture all shared-template dependence. The population bootstrap captures variation across simulated seeds on a fixed suite. Neither interval includes task-design uncertainty, price changes, billing errors, or real model variation. Larger simulated sample counts cannot remove those gaps.
\paragraph{Cost is not value.} A verified task can still have little business value. Conversely, a failed exact-state check can accompany partially useful work. The ratio does not measure user satisfaction, harm, or opportunity cost. Applications should specify acceptable outcomes before deciding how much verified success is worth.
\paragraph{Assistance and review.} An AI coding assistant helps prepare this manuscript and its experiment scripts. The study has no independent human annotation or external replication. The author remains responsible for the claims. The available traces enable readers to inspect the evidence rather than rely on that preparation process.
\section{Citation and availability}
Code, study plan, generated data, and reproduction instructions are available at \url{https://github.com/pooriaarab/benchmark-cost-per-task}. The paper and generated figures use CC-BY-4.0; benchmark code and tasks retain their existing MIT license. This is a preprint packet, not a venue acceptance or an archived publication. Cite the title and author shown here; no new publication DOI is asserted.\par
\section{Authors}
Pooria Arab is an independent researcher, ORCID \AuthorOrcid{}. The author maintains the benchmark and defines the present study. That role is a potential source of design bias. Independent replications should preserve the frozen tasks and report deviations before interpreting differences.\par
\bibliographystyle{plainnat}
\bibliography{references}
\end{document}
